\subsection{群的局部同态的扩张}

定义 1. --- 给定一个拓扑群 G 与一个群 G'，G 到 G' 中的局部同态是从 G 的单位元的一个邻域 V 到 G' 中的一个映射 f，使得对 V 中满足 xy ∈ V 的每对点 x、y，有 f(xy) = f(x)f(y)。

若 G' 是拓扑群且 f 是连续映射，则称 f 是连续局部同态（或拓扑群的局部态射）。

若 G 与 G' 是拓扑群，从 G 到 G' 的局部同构的概念已在 TG, III, p. 6, 定义 2 中定义。G 到 G' 的一个局部同构是从 G 的单位元的一个邻域 V 到 G' 的单位元的一个邻域 V' 上的一个同胚 f，使得 f 与 f 的逆映射都是局部同态。

命题 1. --- 设 G 与 G' 是拓扑群，p: G → G' 是群同态。为使 p 使 G 成为 G' 的一个覆叠，必要且充分的是：p 在 G 的单位元的某个适当邻域上的限制是 G 到 G' 的一个局部同构。

该条件由 TG, III, p. 6 的命题 3 是必要的。反之，假设 \(p\) 诱导一个从 G 的单位元 \(e\) 的开邻域 V 到 G' 的单位元的一个邻域 V' 上的同胚。于是映射 \(p\) 连续（TG, III, p. 15, 命题 23）且开（TG, III, p. 16, 命题 24）。\(p\) 的像 H 是 G′ 的一个包含 V 的子群，故在 G′ 中既开又闭（TG, III, p. 7, 推论）。设 N 是 \(p\) 的核；我们有 N∩V = \{e\}；因此 N 离散（同前， 命题 5）。于是，\(p\) 使 G 成为 H 的一个覆叠，以群 N 为主（I, p. 100, 定理 1 的推论 3）。由于 H 在 G' 中既开又闭，G'-空间 (G, p) 是一个覆叠。

命题 2. --- 设 G 是连通拓扑群，G' 是群，\(f: V \to G'\) 是 G 到 G' 中的局部同态，其中 V 是 G 的单位元的一个连通邻域。则存在一个连通拓扑群 H、一个拓扑群的态射 \(p\colon\mathrm{H}\to\mathrm{G}\)（使 \((\mathrm{H},p)\) 是 G 的一个覆叠）以及一个群同态 \(\varphi\colon\mathrm{H}\to\mathrm{G}'\)，使得由 \(y\in p^{-1}(\mathrm{V})\)（满足 \(f(p(y))=\varphi(y)\)）所成的集合是 H 中单位元的一个邻域。

若 G' 是拓扑群且映射 f 连续，则这样的同态 \(\varphi\) 是连续的。

引理 1. — 设 G 是拓扑群，V 是 G 的单位元 e 的一个连通邻域。对 e 在 G 中的每个邻域 U 及每个 \(x \in V\)，存在一个整数 \(n \in N\) 及 U 中的元素 \(u_{1}, \ldots, u_{n}\)，使 \(u_{1} \ldots u_{n} = x\)，且 \(u_{1} \ldots u_{k} \in V\)（对每个满足 \(1 \leqslant k \leqslant n\) 的整数 k）。

我们可以假设 U 开且包含于 V。设 A 表示满足引理条件的 \(x \in V\) 所成的集合。若 \(x \in A\) 且 \(y \in xU \cap V\)，则 \(y \in A\)，由此得 \(AU \cap V \subset A\)；这表明 A 在 V 中开。设 \(x \in V\) 使 \(xU^{-1} \cap A \neq \varnothing\)；于是有 \(x \in AU \cap V\)，因此 \(x \in A\)。于是，若 \(x \in V\) 且 \(x \notin A\)，则 \(xU^{-1} \cap A = \varnothing\)，故 A 在 V 中闭。由于 \(e \in A\) 且 V 连通，得出 A = V。

现在来证明命题。记 \(j\) 为映射 \(g \mapsto (g, f(g))\)（从 V 到 \(\mathbf{G} \times \mathbf{G}'\)），并设 H 是 \(\mathbf{G} \times \mathbf{G}'\) 的一个子群（由 \(j(\mathrm{V})\) 生成）。

设 U 是 \(e\) 在 G 中的一个邻域，包含于 V。设 \(x \in V\)；由引理 1，存在 \(u_1, \ldots, u_n \in U\)，使 \(x = u_1 \ldots u_n\)，且 \(u_1 \ldots u_k \in V\)（对每个整数 \(k\)，满足 \(1 \leqslant k \leqslant n\)）。由归纳法，有 \(f(u_1 \ldots u_k) = f(u_1) \ldots f(u_k)\)（对每个整数 \(k\)，\(1 \leqslant k \leqslant n\)）。特别地，\(j(x)\) 属于由 \(j(U)\) 生成的子群。由此得出 H 由 \(j(U)\) 生成。

设 \(\mathcal{B}\) 是 H 的形如 \(j(\mathrm{U})\) 的子集所成的集合，其中 U 是 \(e\) 在 V 中的一个邻域。我们来证明 H 上存在唯一的拓扑，它与 H 的群结构相容，且使 \(\mathcal{B}\) 是单位元的邻域滤子的基。为此，只需证明集合 \(\mathcal{B}\) 满足 III, p. 4 中的条件 \((\mathrm{GV}_{\mathrm{I}}^{\prime})\)、\((\mathrm{GV}_{\mathrm{II}}^{\prime})\) 与 \((\mathrm{GV}_{\mathrm{III}}^{\prime})\)。

设 U 是 e 在 G 中的一个邻域，包含于 V；存在 e 的一个邻域 \(U'\)，使 \(U' \cdot U' \subset U\)。对 \(U'\) 的每对点 x、y，有 \(xy \in V\) 且 \(f(xy) = f(x)f(y)\)。于是，\(j(\mathrm{U}') \in \mathcal{B}\) 且 \(j(\mathrm{U}') \cdot j(\mathrm{U}') \subset j(\mathrm{U})\)。这表明条件 \((\mathrm{GV}_{\mathrm{I}}')\) 满足。

于是集合 \(U'' = V \cap U^{-1}\) 是 e 在 V 中的一个邻域，且对 \(x \in U''\)，有 \(x^{-1} \in U\) 及 \(f(x^{-1}) = f(x)^{-1}\)。因而 \(j(\mathrm{U}'')^{-1} \subset j(\mathrm{U})\)，这表明条件 \((\mathrm{GV}_{\mathrm{II}}' )\) 成立。

最后，取定 e 在 V 中的一个邻域 U，使 \(U = U^{-1}\) 且 \(U^{3} \subset V\)。设 W 是 e 在 U 中的一个邻域，设 \(h = (g, f(g))\) 是 \(j(\mathrm{U})\) 的一个元素。存在 e 的一个包含于 W 的邻域 \(W'\)，使 \(gW'g^{-1} \subset W\)。于是，\(j(\mathrm{W}') \in \mathcal{B}\) 且 \(hj(\mathrm{W}')h^{-1} \subset j(\mathrm{W})\)，因为有 \(f(gxg^{-1}) = f(g)f(x)f(g^{-1})\)（对 \(x \in W'\)）。

设 \(h \in H\)。由于 U 是 e 的包含于 V 的邻域，\(j(\mathrm{U})\) 生成 H；由于 U 对称，存在 U 中的元素 \(u_{1}, \ldots, u_{n}\)，使 \(h = j(u_{1}) \ldots j(u_{n})\)。对 n 用归纳法，存在 e 的一个包含于 W 的邻域 \(W'\)，使 \(hj(\mathrm{W}')h^{-1} \subset j(\mathrm{W})\)。于是条件 \((\mathrm{GV}_{\mathrm{III}}')\) 满足。

然后给群 H 赋以这个拓扑。

记 \(p \colon \mathrm{H} \to \mathrm{G}\) 为第一个投影 \(\mathrm{G} \times \mathrm{G}' \to \mathrm{G}\) 在 H 上的限制。这是一个群同态。对 \(e\) 的每个包含于 V 的邻域 U，集合 \(p^{-1}(\mathrm{U})\) 包含 H 的单位元的邻域 \(j(\mathrm{U})\)，故 \(p\) 是拓扑群的一个连续同态。对 \(e\) 的每个包含于 V 的邻域 U，有 \(p(j(\mathrm{U})) = \mathrm{U}\)，因此 \(p\) 是开映射。由于 G 连通，\(p\) 是满射的。它的核是离散的，因为它除单位元外与 \(j(\mathrm{V})\) 不相交。于是由推论 3（I, p. 100）得出，G-空间 \((\mathrm{H}, p)\) 是一个覆叠。

设 \(\varphi\) 是第二个投影 \(G \times G' \to G'\) 在 H 上的限制。若 \(g \in V\)，有 \((g, f(g)) = j(g) \in j(V)\) 且 \(\varphi(g, f(g)) = f(g)\)，因而有 \(\varphi(y) = f(p(y))\)（对 \(y \in j(V)\)）。

进而假设 G' 是拓扑群且映射 f 在 e 处连续。于是同态 \(\varphi\) 在 H 的单位元处连续，因而连续（TG, III, p. 15, 命题 23）。

推论 1. --- 设 G 是单连通拓扑群，G' 是群。设 V 是 G 中单位元的一个连通邻域，f: V → G' 是一个局部同态。存在唯一的群同态 h: G → G' 延拓 f。若 G' 是拓扑群且映射 f 连续，则同态 h 连续。

群 G 连通。设 H、\(\varphi\) 与 p 如命题 2 中所示。由于 G 单连通，H 是 G 的一个可平凡化覆叠（I, p. 124, 定义 3）；由于 H 连通且非空，映射 \(p\) 是拓扑群的一个同构。置 \(h = \varphi \circ p^{-1}\)；映射 \(h\) 是群同态，且存在 G 的单位元 \(e\) 的一个包含于 V 的开邻域 U，使 \(f|U = h|U\)。由引理 1 得出，\(f\) 与 \(h\) 在 V 上一致。换句话说，映射 \(h\) 延拓映射 \(f\)。

若 G' 是拓扑群且映射 f 连续，则同态 \(\varphi\) 连续，故 h 也连续。

我们来证明这种延拓的唯一性。G 中使从 G 到 G′ 的两个同态一致的点所成的集合是 G 的一个子群。由于群 G 连通，G 的每个包含单位元的一个邻域的子群都等于 G（TG, III, p. 8, 命题 6）。h 的唯一性由此得出。

注 1. --- 当 G = R 时，这个推论由 TG, V, p. 3 的命题 6 得出。

推论 2. --- 两个局部连通且单连通的拓扑群若是局部同构的，则是同构的。

设 G 与 G' 是局部连通拓扑群，V 与 V' 分别是 G 与 G' 的单位元的连通邻域，\(f: V \to V'\) 是一个同胚，且是 G 到 G' 的一个局部同构。由推论 1，存在唯一的连续群同态 \(\varphi: G \to G'\) 延拓 f，且存在唯一的连续群同态 \(\varphi': G' \to G\) 延拓 \(f^{-1}\)。同态 \(\varphi' \circ \varphi\) 与 \(\operatorname{Id}_{\mathrm{G}}\) 在 e 于 G 中的一个邻域上一致，故相等，因为 G 连通。同样，\(\varphi \circ \varphi' = \operatorname{Id}_{\mathrm{G}'}\)，推理由此得证。

推论 3. --- 设 G 是单连通拓扑群，V 是 G 的单位元的一个连通邻域。以集合 V 为生成集，并以关系元集 \(\mathbf{r}\) 为诸 \(xyz^{-1}\)（其中 \((x,y,z)\) 取遍 V 中满足 \(xy = z\) 的元素三元组）所成的族，便定义了 G 的一个表现。

设 \(\mathrm{F}(\mathrm{V}, \mathbf{r})\) 是自由群 \(\mathrm{F}(\mathrm{V})\) 关于包含诸元素 \(xyz^{-1}\)（其中 \((x, y, z) \in \mathrm{V} \times \mathrm{V} \times \mathrm{V}\) 且 xy = z）的最小正规子群的商群（A, I, p. 86）。记 \(f: V \to F(V, r)\) 与 \(g: F(V, r) \to G\) 为典范映射。按构造，映射 f 是 G 到 \(\mathrm{F}(V, \mathbf{r})\) 的一个局部同态。由推论 1，存在唯一的群同态 \(\overline{f}: G \to F(V, r)\) 延拓 f。

由于群 \(F(V, \mathbf{r})\) 由 \(f(V)\) 生成，同态 \(\overline{f}\) 是满射的。对 \(x \in V\)，有 \(g(\overline{f}(x)) = g(f(x)) = x\)；由于 V 生成 G，\(g \circ \overline{f} = \operatorname{Id}_{\operatorname{G}}\)，这证明 \(\overline{f}\) 是单射的。因此它是一个同构。
% semantic-fragments: legacy-input-seam
\relax{}
